\section{Go, Bitter Lesson y la Inference Compute general}

El go ofrece un segundo punto de comparación. Es un juego de información perfecta, sin cartas ocultas, y aun así depende de search para convertir la policy/value network en acciones sólidas. El ponente se apoya en la ablation de AlphaGo Zero y en la Bitter Lesson para ampliar la «búsqueda», que deja de ser un algoritmo específico de juegos y adopta una definición más amplia: seguir invirtiendo cómputo después de recibir la entrada para mejorar la salida de esa ocasión.

\subsection{AlphaGo Zero: la red no sustituye a la búsqueda}

Las barras de la figura muestran la fuerza de juego de cada versión del sistema. Full AlphaGo Zero usa al mismo tiempo la red entrenada y Monte Carlo Tree Search. Al retirar MCTS, el desempeño cae drásticamente por debajo del nivel expert-human, aunque la red provenga de un self-play sólido. Esto no significa que MCTS pueda trasladarse directamente a cualquier tarea, sino que las heurísticas obtenidas en el entrenamiento y el despliegue durante la inferencia se complementan.

\lecturefigure{slide-06-search-go.jpg}{Ablación de búsqueda de AlphaGo Zero: la diferencia entre el Full system y la versión no-search es marcada}{Diapositiva V006 recuperada de la grabación oficial; Silver et al., Science 2017}

\begin{knowledgebox}{Lectura de la figura: la altura de la barra es la fuerza de juego, no la loss de entrenamiento}
Esta figura compara el nivel final en partidas. La policy/value network proporciona el move prior y la leaf evaluation, y MCTS los combina en decisiones más confiables sobre cada posición. Sin búsqueda, el modelo no queda «sin saber nada», pero no logra organizar sus predicciones locales en secuencias de acciones de largo alcance lo bastante sólidas. Por eso, los parámetros del modelo y las search visits no pueden tratarse como una sola escala totalmente intercambiable.
\end{knowledgebox}

\subsection{Cómo entender la analogía de 100,000 veces de la clase}

El ponente usa una regla muy aproximada para dar una intuición del orden de magnitud: si cada aumento de diez veces en cierta dimensión del base system aporta solo unos 50 Elo, mientras que search aporta unos 250 Elo, entonces para lograr la misma diferencia de Elo se necesitan cinco aumentos de diez veces, es decir,
\begin{equation}
R_{\text{equiv}}
\approx 10^{\Delta E/50}
=10^{250/50}=10^5.
\end{equation}
$\Delta E$ es la ganancia de Elo que se quiere sustituir, 50 es la pendiente empírica citada en clase y $R_{\text{equiv}}$ es el factor de escala aproximado. El valor de este cálculo es recordarnos que no debemos ignorar un desplazamiento grande de la curva, no afirmar que search equivale siempre a cien mil veces el cómputo de entrenamiento en cualquier tarea.

\teachervoice{El docente presenta explícitamente «unas 100,000 veces» como una extrapolation útil para la discusión, no como una ley. La tesis real es esta: si solo optimizamos a lo largo del eje de base-model scaling, podemos pasar por alto un inference axis mucho más barato, pero que exige un nuevo diseño de sistemas.}

\subsection{Bitter Lesson: métodos generales y cómputo escalable}

La Bitter Lesson de Richard Sutton sostiene que, a largo plazo, el progreso más confiable en IA suele provenir de métodos generales capaces de aprovechar un cómputo en crecimiento constante, en particular learning y search, y no de codificar de forma rígida grandes cantidades de conocimiento humano en el sistema. La clase la cita para situar el póquer y el go en una misma línea histórica: el conocimiento del dominio sigue siendo útil, pero debe estar al servicio de un proceso de cómputo escalable, no reemplazarlo.

\lecturefigure{slide-07-bitter-lesson.jpg}{The Bitter Lesson: search y learning son métodos generales que escalan con el cómputo}{Diapositiva V007 recuperada de la grabación oficial; Richard Sutton, 2019}

\begin{warningbox}{«Método general» no equivale a «sin estructura del dominio»}
Libratus, Pluribus, AlphaGo Zero y Cicero incluyen una gran cantidad de modelado del dominio. En esta clase, la Bitter Lesson sirve para elegir la dirección de optimización: dar prioridad a mecanismos capaces de aprovechar más cómputo, no afirmar que toda hand-crafted representation debe eliminarse de inmediato.
\end{warningbox}

\subsection{El siguiente objetivo: ¿es posible la general inference computation?}

La octava slide lleva la pregunta directamente a los modelos de lenguaje: si el base model ya es muy grande, ¿puede seguir mejorándose cada salida con un inference compute mucho más barato que el entrenamiento? Y si una respuesta tiene un valor muy alto, ¿debería permitirse que el sistema dedique minutos, horas o incluso más tiempo? Aquí \term{search} es un término amplio que no se limita a estructuras de árbol explícitas: también puede consistir en generar, verificar y corregir repetidamente, o en recurrir a la retroalimentación del entorno.

\lecturefigure{slide-08-next-goal-generality.jpg}{De la búsqueda en juegos a la inference-time computation general}{Diapositiva V008 recuperada de la grabación oficial; Stanford CS25 Lecture 14}

\begin{importantbox}{Presupuesto de inferencia sensible al valor}
Un inference budget razonable no debería depender solo de la longitud de la entrada, sino también del costo del error, de la posibilidad de verificar el resultado y del costo de oportunidad. Recomendar una canción admite una respuesta rápida; el diseño de fármacos, el enrutamiento de chips o una negociación crucial justifican generar más candidatos y evaluarlos con mayor rigor. El objetivo del sistema es maximizar «el valor marginal que aporta el cómputo adicional», no hacer más lentas todas las solicitudes por igual.
\end{importantbox}

\teachervoice{El ponente considera chain-of-thought una forma de cómputo adicional muy general, pero todavía limitada: el modelo desarrolla el problema generando tokens intermedios, pero no necesariamente cuenta con un planificador que permita retroceder, verificar o mantener un estado a largo plazo. La pregunta abierta de la clase es si existe un planning mechanism más general y más confiable que CoT.}

\subsection{Resumen de la sección}

La ablación en go demuestra que una network sólida sigue necesitando search, la Bitter Lesson ofrece la explicación histórica y la generality slide orienta la agenda de investigación hacia sistemas de inferencia con presupuesto variable. A continuación, Diplomacy eleva la dificultad un nivel más: los otros «estados» del entorno son personas que hablan, cooperan y también cambian de opinión.
